% from ega4-8-fr.tex
\oldpage[IV-3]{16}
sont commutatifs~; on en déduit donc par passage à la limite inductive un
diagramme commutatif
\begin{equation}
\xymatrix{
  \varinjlim\mathfrak P(S_\lambda) \ar[r]^{u_{\mathfrak P}}
    \ar[d]_{\varinjlim g_\lambda^{-1}}
    & \mathfrak P(S) \ar[d]^{g^{-1}} \\
  \varinjlim\mathfrak P(X_\lambda) \ar[r]_{v_{\mathfrak P}}
    & \mathfrak P(X)
}
\tag{8.3.10.1}\label{IV.8.3.10.1-fr}
\end{equation}
et il résulte de (1.8.2) que l'on a des diagrammes analogues en remplaçant
$\mathfrak P$ par $\mathfrak C$, $\mathfrak D_c$, $\mathfrak F_c$ ou
$\mathfrak{LF}_c$.
\end{env}

Il résulte de (8.3.5) que sous l'hypothèse que pour un $\alpha\in L$,
$S_\alpha$ est \emph{quasi-compact}, l'application canonique (8.3.9.2) est
\emph{injective}. En outre~:

\begin{theorem}[8.3.11]
\label{IV.8.3.11-fr}
Supposons qu'il existe un $\alpha\in L$ tel que $S_\alpha$ soit quasi-compact
et quasi-séparé. Alors les applications canoniques (8.3.9.2), (8.3.9.3),
(8.3.9.4) et (8.3.9.5) sont \emph{bijectives}.
\end{theorem}

En vertu de la remarque précédente, il reste à prouver que ces applications
sont \emph{surjectives}~; comme toute partie constructible de $S$ est réunion
finie d'ensembles de la forme $U\cap\complement V$, où $U$ et $V$ sont
ouverts et constructibles, il suffira de prouver que (8.3.9.3) est surjective
pour qu'il en soit de même de (8.3.9.2) (et aussi de (8.3.9.4), par passage
aux complémentaires). Or, comme les morphismes $S_\lambda\to S_\alpha$ et
$S\to S_\alpha$ sont affines, donc séparés, les $S_\lambda$ pour
$\lambda\geq\alpha$ et $S$ sont quasi-compacts et quasi-séparés (1.2.2), et
l'on sait que les parties ouvertes constructibles dans un tel préschéma ne
sont autres que les parties ouvertes quasi-compactes (1.8.1). La conclusion
résulte donc de (8.2.11) sauf pour l'application (8.3.9.5). Pour prouver que
cette dernière est surjective, considérons une partie $Z$ localement fermée
et constructible dans $S$~; $Z$ est donc quasi-compact
$(0_{\mathrm{III}},9.1.4)$. Comme tout point $x\in Z$ admet par hypothèse un
voisinage ouvert quasi-compact $V_x$ dans $S$ tel que $Z\cap V_x$ soit fermé
dans $V_x$, on peut couvrir $Z$ par un nombre fini des $V_x$~; autrement dit,
il y a un ouvert quasi-compact $U$ contenant $Z$ et tel que $Z$ soit fermé
dans $U$~; comme $Z$ est constructible dans $S$, il l'est aussi dans $U$
$(0_{\mathrm{III}},9.1.8)$. On sait (8.2.11) qu'il y a un indice $\lambda$ et
un ouvert quasi-compact $U_\lambda$ dans $S_\lambda$ tels que
$U=u_\lambda^{-1}(U_\lambda)$. Appliquant à $U$ (qui est limite projective
des $U_\mu=u_{\lambda\mu}^{-1}(U_\lambda)$ pour $\mu\geq\lambda$) le fait
que l'application (8.3.9.4) est surjective, on voit qu'il existe un
$\mu\geq\lambda$ et un ensemble fermé constructible $Z_\mu$ dans $U_\mu$
tels que $Z=u_\mu^{-1}(Z_\mu)$. Mais comme l'immersion canonique
$U_\mu\to S_\mu$ est quasi-compacte par hypothèse (1.2.7), elle est de
présentation finie (1.6.2, (i)), et $Z_\mu$ est aussi une partie constructible
\emph{de} $S_\mu$ en vertu de (1.8.4) et (1.8.1)~; comme $Z_\mu$ est
évidemment localement fermée dans $S_\mu$, cela achève la démonstration.

\begin{corollary}[8.3.12]
\label{IV.8.3.12-fr}
Supposons qu'il existe un $\alpha$ tel que $S_\alpha$ soit quasi-compact, et
soit, pour tout $\lambda$, $Z_\lambda$ une partie constructible de
$S_\lambda$ telle que
$Z_\mu=u_{\lambda\mu}^{-1}(Z_\lambda)$ pour $\mu\geq\lambda$. Si
$Z=u_\lambda^{-1}(Z_\lambda)$,
